%%=============================================================================
%% De inhoud van je projectvoorstel
%%=============================================================================
%%
%% Zeven rubrieken, samen hoogstens twee bladzijden. Houd het kort: per
%% rubriek volstaan één tot drie zinnen. Een voorstel is geen rapport - het
%% moet je begeleider in vijf minuten laten beslissen of dit een geschikt
%% graduaatsproject is.
%%
%% Stem je voorstel af met je mentor VOORDAT je het indient. Je mentor
%% bevestigt dat dit is wat het bedrijf bedoelt en dat de tien werkdagen
%% haalbaar zijn.

%---------- 1. Context en probleemstelling ------------------------------------

\section{Context en probleemstelling}%
\label{vst:context}

% Waar loop je stage, wat doet dat team, en welk concreet probleem of welke
% concrete behoefte is de aanleiding? Noem de mensen voor wie het een probleem
% is. Vermijd algemeenheden als "bedrijven hebben nood aan".

Ik loop stage bij Attentia en ben lid van het AI-team. Dit team houdt in dat ze systemen binnen attentia willen aanpassen, vernieuwen en optimaliseren door AI-technieken te implementeren in de systemen die al aanwezig zijn. Het gaat hier dan zowel over interne programma's en tools als externe, bijvoorbeeld voor de HR-klanten of wellbeing klanten.

%---------- 2. Doelstelling ---------------------------------------------------

\section{Doelstelling}%
\label{vst:doelstelling}

% Wat bestaat er op het einde van dit project dat er vandaag niet is? Formuleer
% het als een resultaat waarvan iemand kan vaststellen dat het er is, en
% vermeld waaraan je zal merken dat het geslaagd is.

Op het einde van het project bestaat er een self-service oplossing waarmee een klant zelf, via DOTS of een platform, een tijdelijk FTP-account aanvraagt. De backend-setup gebeurt daarna automatisch: het account wordt aangemaakt, de toegang wordt geconfigureerd en de inloggegevens worden veilig aan de klant bezorgd via mail. Na een vastgelegde periode (20 dagen) worden de bestanden automatisch verwijdert en na lange inactiviteit het account ook. Het project is geslaagd als een volledige uitwisseling (aanvraag, account, upload of download en opruiming) kan verlopen zonder dat iemand van IT tussenbeide komt, en als dat in een end-to-end test aantoonbaar werkt.

%---------- 3. Technologieën --------------------------------------------------

\section{Technologieën}%
\label{vst:technologieen}

% Met welke talen, frameworks, databanken en diensten ga je werken? Welke
% liggen vast door de omgeving van het bedrijf, en welke kies je zelf? Vermeld
% ook wat nieuw voor je is: technologische verbreding wordt gewaardeerd.
%
% Vermeld hier ook of je AI denkt in te zetten, en waarvoor. Dat mag en wordt
% zelfs verwacht; in je rapport verantwoord je later wat je precies gedaan hebt.
%
% Verwerkt je toepassing persoonsgegevens, of zit er AI in wat je bouwt? Zeg
% dat dan nu al in één zin: dan weet je begeleider meteen dat de AVG of de
% AI Act meespeelt, en kan dat mee bekeken worden bij de goedkeuring.

WingFTP is de FTP-server en ligt vast door de omgeving van het bedrijf, net als DOTS als vertrekpunt voor de aanvraag. De backend zal in C\# geschreven zijn en voor de communicatie zorgen tussen de frontend en de FTP-server. De frontend wordt gemaakt in react tijdens het prototyping en later vermoedelijk met Angular, aangezien de andere devs hiermee werken. De scripting en automatisatie van het accountbeheer en de koppeling tussen aanvraag en FTP-server doe ik met powershell en vermoedelijk n8n (nog niet 100 procent zeker). Daarbij komen netwerkdesign en firewallconfiguratie, en serverbeheer van de onderliggende infrastructuur. De oplossing zelf bevat geen AI wordt mee bekeken met mijn begeleider.

%---------- 4. Afbakening -----------------------------------------------------

\section{Afbakening}%
\label{vst:afbakening}

% DE BELANGRIJKSTE RUBRIEK. Je stage en je graduaatsproject lopen bij hetzelfde
% bedrijf; dit is waar je ze uit elkaar houdt.
%
% Beschrijf:
%   - wat je dagelijkse stagewerk is;
%   - wat het afgebakende geheel is dat JIJ zelfstandig analyseert, ontwerpt
%     en oplevert - dat is je graduaatsproject;
%   - wat er uitdrukkelijk buiten valt.
%
% Toets voor jezelf: welke analyse en welk ontwerp maak jij zelf? Kan je daar
% niets over schrijven, dan heb je een taak en geen project.

Mijn stage bestaat uit meerdere projecten, waaronder het Zero Touch-project, waarvan de FTP-oplossing deel uitmaakt, en de ontwikkeling van een AI-agent. Mijn graduaatsproject is de self-service FTP-oplossing als afgebakend geheel, end-to-end. Dat gaat verder dan code schrijven. Ik analyseer zelf de noden van klanten en IT, onderzoek de mogelijke aanpak en maak het ontwerp: de flow van aanvraag tot opruiming, de levenscyclus en beveiliging van accounts, de serverinfrastructuur, het netwerk en de firewallregels, en de communicatie met de klant. Daarna bouw ik de oplossing, test ik ze en documenteer ik ze.

%---------- 5. Verwachte meerwaarde -------------------------------------------

\section{Verwachte meerwaarde}%
\label{vst:meerwaarde}

% Wat heeft het bedrijf hieraan? Wat blijft er van jouw werk over wanneer je
% stage afgelopen is? Denk aan tijdwinst, minder fouten, nieuwe mogelijkheden,
% lagere kosten, of een basis waarop verder gebouwd kan worden.

oor IT verdwijnt het handmatig aanmaken en opruimen van accounts bij eenmalige uitwisselingen, wat veel tijd bespaart. Klanten krijgen sneller toegang, zonder te wachten op IT. Doordat accounts na een vaste periode automatisch worden uitgeschakeld en verwijderd, blijven er geen vergeten accounts met toegang open staan, wat de veiligheid verhoogt. Na mijn stage blijven de werkende oplossing, de documentatie en de beheerdershandleiding over. Die vormen een basis om het Zero Touch-principe, waarbij klanten zelf zaken regelen zonder menselijke tussenkomst, verder uit te breiden.

%---------- 6. Planning en werkbelasting --------------------------------------

% Het graduaatsproject is ongeveer tien werkdagen, samen ongeveer 100 uren.
% Verdeel die over de fasen van je project en geef per fase aan wat je oplevert.
% Vermeld ook of je die dagen gefocust opneemt of gespreid over de stageperiode,
% en dat je dit met je mentor afgesproken hebt.

\newcommand{\vstplanningtabel}{%
  \centering
  \begin{tabular}{l p{0.6\textwidth} r}
    \toprule
    \textbf{Fase} & \textbf{Wat lever je op} & \textbf{Uren} \\
    \midrule
    Analyse en vereisten      & overzicht van de noden van klanten en IT en de vereisten van de oplossing & 12 \\
    Technologiekeuze          & verantwoorde keuze van scripttaal, koppeling met de aanvraag en manier van veilige bezorging van inloggegevens & 8 \\
    Ontwerp                   & technisch ontwerp: aanvraagflow, accountlevenscyclus, beveiliging, netwerk en firewallregels & 16 \\
    Realisatie                & werkende oplossing: geautomatiseerd accountbeheer, koppeling met de aanvraag, serverconfiguratie & 38 \\
    Testen en validatie       & testverslag van de end-to-end test en de controle van beveiliging en opruiming & 10 \\
    Rapport, poster en demo   & eindrapport, technische documentatie, beheerdershandleiding, poster en demo & 16 \\
    \midrule
    \textbf{Totaal}           &                              & \textbf{100} \\
    \bottomrule
  \end{tabular}
  \caption[Planning en werkbelasting]{\label{tab:vst-planning}%
    Vervang deze verdeling door je eigen planning.}}

% In het tweekolommige voorstel past de tabel niet in één kolom. Daarom
% beginnen titel en tabel samen bovenaan een nieuwe pagina, over de volle
% breedte, net zoals de samenvatting. In het rapport (één kolom) blijft het
% een gewone tabel.
\makeatletter
\if@twocolumn
  \twocolumn[{%
    \section{Planning en werkbelasting}\label{vst:planning}%
    \def\@captype{table}%
    \vstplanningtabel
    \bigskip}]
\else
  \section{Planning en werkbelasting}%
  \label{vst:planning}
  \begin{table}[htbp]
    \vstplanningtabel
  \end{table}
\fi
\makeatother

%---------- 7. Risico's -------------------------------------------------------

\section{Risico's}%
\label{vst:risicos}

% Wat kan dit project doen mislukken, en wat doe je dan? Denk aan: toegang tot
% gegevens of systemen die je nodig hebt, afhankelijkheid van een collega of
% een derde partij, technologie die je nog moet leren, of drukte op de
% werkvloer waardoor je dagen niet vrijkomen.
%
% Twee tot drie risico's met telkens één zin over wat je terugvalplan is.

Een eerste risico is dat WingFTP minder mogelijkheden voor automatisch accountbeheer biedt dan verwacht, terwijl ik er zelf nog mee moet leren werken. Als dat zo blijkt, onderzoek ik dit vroeg in de analyse- en keuzefase en val ik terug op automatisatie via scripting of configuratiebestanden in plaats van een rechtstreekse koppeling.

Een tweede risico is dat ik voor de koppeling met DOTS, de serverconfiguratie en de firewallregels afhankelijk ben van toegang en medewerking van collega's van IT en security. Ik regel die toegang en afspraken in het begin van het project met mijn mentor en werk in de tussentijd verder in een testomgeving.

Een derde risico is de beperkte tijd: met ongeveer 100 uren is het end-to-end karakter van het project een uitdaging, zeker als dagen door drukte op de werkvloer niet vrijkomen. Daarom bouw ik eerst een minimale werkende versie (aanvraag, automatisch account en automatische verwijdering) en pak ik uitbreidingen pas daarna aan, zodat er altijd een bruikbaar resultaat is.
